\section{Introduction}

Aging infrastructure worldwide---including bridges, industrial plants, tunnels, and civil constructions---requires frequent inspection and maintenance to prevent catastrophic failures. The manual execution of these tasks exposes human workers to hazardous conditions, incurs high operational costs, and is limited in coverage and repeatability \citep{Halder2023Construction,Wang2025Autonomous}. Autonomous robotic inspection has therefore emerged as a critical technology, offering scalability, continuous operation, and the capacity to acquire high-precision data in environments that are inaccessible or dangerous for humans \citep{Halder2025Crack}. The integration of autonomous sensing platforms into non-destructive testing (NDT) workflows has been identified as a key enabler for next-generation infrastructure monitoring, combining visual detection, inertial navigation, and point-cloud mapping into a unified inspection pipeline.

Among existing robotic platforms, legged robots (and quadrupeds in particular) have proven the most effective for semi-structured environments that combine flat corridors with obstacles, stairs, and uneven ground. ANYmal, developed at ETH Z\"{u}rich, was specifically designed for harsh industrial environments and has demonstrated reliable operation in oil and gas plants and search-and-rescue scenarios \citep{Hutter2017ANYmal}. More recently, Team CERBERUS deployed a multi-agent system of ANYmal robots augmented with LiDAR, cameras, and IMU to win the DARPA Subterranean Challenge, establishing the quadruped as the reference platform for autonomous exploration in confined and semi-structured spaces \citep{Tranzatto2022CERBERUS}. Despite these advances, quadrupedal robots face a fundamental efficiency limitation: walking gaits are energetically costly on flat surfaces, which restricts their operational autonomy for tasks that combine open traversal (e.g., factory floors) with the negotiation of obstacles or slopes \citep{Vogel2024Ladder}. This trade-off between terrain adaptability and locomotion efficiency motivates the development of hybrid platforms capable of switching between legged and wheeled modes according to terrain properties.

Wheeled-legged robots address this limitation by integrating actuated limbs (providing the dexterity required to negotiate irregular or sloped terrain) with wheels that constitute the most efficient effector on flat surfaces \citep{Bjelonic2023Survey}. The ANYmal-on-wheels family, developed at ETH Z\"{u}rich, is the most representative line in this category. \citet{Bjelonic2019KeepRollin} introduced the first whole-body controller capable of simultaneously combining walking and rolling with non-steerable torque-controlled wheels; subsequent work \citep{bjelonic2019} extended this to online trajectory optimization for hybrid maneuvers validated in the DARPA SubT, and \citet{Bjelonic2022OfflineMotion} later incorporated offline motion libraries and MPC to achieve complex skills such as stair climbing and dynamic 180° turns. The most recent and impactful result in this family, \citet{Lee2024Science}, demonstrated a complete pipeline of adaptive locomotion control, mobility-aware local planning, and city-scale global navigation via deep reinforcement learning (RL) with privileged learning, achieving smooth walking-to-driving transitions as an emergent property of a single learned policy. The CENTAURO platform from the Italian Institute of Technology \citep{Kashiri2019CENTAURO,Klamt2019CENTAURO} ---a centaur-type quadruped with steerable omnidirectional wheels and a high-payload manipulator arm--- offers a complementary design point with explicit disaster-response validation and provides the best published quantitative evidence for the energy/terrain trade-off that motivates a hybrid robot. At the compact end of the spectrum, Ascento \citep{Klemm2019Ascento} is a bipedal wheeled-legged platform for indoor inspection, demonstrating that the design space spans from bipedal (Ascento) to centaur (CENTAURO) and quadrupedal (ANYmal-wheels) morphologies. Recent platforms such as those published by \citet{Liu2024Max},  \citet{Cui2025FLORES}, and \citet{Jiang2025RTaichi} further expand the design space with alternative articulation schemes and RL-based controllers, confirming that the field is converging on hybrid locomotion as the dominant paradigm for versatile robot operation.

A critical observation emerges from this body of work: in all published wheeled-legged systems, locomotion mode switching is resolved either heuristically through hand-coded threshold logic \citep{Bjelonic2019KeepRollin,Kashiri2019CENTAURO}, by model predictive control \citep{bjelonic2019,Bjelonic2022OfflineMotion}, or by end-to-end RL policies that encode switching implicitly as a black-box behavior \citep{Lee2024Science,Humphreys2025BioInspired}. None of these approaches employs a neurobiologically interpretable arbitration architecture. \revblue{The basal ganglia, a subcortical neural circuit present in all vertebrates, are the canonical neural substrate for action selection: through a salience-driven winner-take-all (WTA) mechanism implemented via disinhibition, they select and sustain a single behavioral channel from a set of competing alternatives while suppressing all others.} To the best of the authors' knowledge, no prior work has applied a basal-ganglia-type architecture to arbitrate among rolling, walking, and omnidirectional locomotion modes in a hybrid quadruped. This constitutes the principal novelty of the present contribution.

Neural control architectures seek to emulate the information-processing structures found in biological systems in order to generate adaptive and robust motor responses \citep{mead1990neuromorphic,ijspeert2008}. The computational model of basal ganglia (BG) ---the GPR model--- proposed by \citet{Gurney2001a,Gurney2001b} reformulates action selection as the resolution of competition among salience-coded channels through a feedforward off-center/on-surround network with selection and control pathways. \citet{Prescott2006Robot} provided the first embodied validation of the GPR model, implementing it in a mobile robot performing a foraging task with five competing behaviors, demonstrating clean switching, appropriate behavioral persistence, and conflict-resolution dynamics analogous to those observed in animals. \citet{Girard2003Basal} compared BG-based action selection against a classical WTA and showed that the BG's thalamocortical loops confer specific adaptive properties such as avoidance of dithering and energetic savings that a simple WTA cannot reproduce. Building upon this theoretical foundation, \citet{sarvestani2013computational} developed a computational network model that operationalizes these selection mechanics. This distinction and its subsequent computational formulation are directly relevant to the present work, as they provide the structural blueprint for our real-time mode arbitration. More recent updates to the GPR paradigm include a biologically parsimonious three-pathway model (Go/NoGo/hyperdirect) \citep{Baston2015Biologically} and a neurorobotic study of simulated dopamine modulation that shows how tonic dopamine levels control the speed and frequency of mode switching \citep{Prescott2024Simulated}, opening a pathway to meta-control of switching aggressiveness in future work. For the rhythmic generation of the quadrupedal walking gait, the present platform integrates central pattern generators (CPGs), neural circuits that produce coordinated rhythmic motor patterns from simple inputs \citep{ijspeert2008}; this combination of BG arbitration above and CPG execution below forms a hierarchical neural stack that is anatomically plausible and computationally modular \citep{Bellegarda2022CPGRL,Manoonpong2013Neural}.

Against this background, the present work proposes a hybrid quadrupedal robotic platform with four three-degree-of-freedom limbs, each incorporating an omnidirectional wheel at the distal segment (tibia), enabling transitions among three discrete locomotion modes: quadrupedal walking (articulated mode), differential rolling, and omnidirectional motion (mobile mode). The platform integrates a sensory suite comprising an IMU, a LiDAR, and a camera for environmental perception. A multilayer perceptron (MLP) processes sensory data to detect terrain instability and compute per-channel salience inputs to the basal-ganglia network, which in turn executes mode selection via a WTA disinhibition mechanism. The complete system is validated in simulation (Gazebo) and on a physical prototype, demonstrating robust autonomous switching across flat surfaces, sloped terrain, and obstacle-laden environments. Relative to the current state of the art, the proposed architecture offers three differentiating properties: \revblue{(i) \emph{biological interpretability}, following directly from the neuroscientific grounding established above;} (ii) \emph{modularity}: perception, arbitration, and motor generation are decoupled, facilitating individual upgrade; and (iii) \emph{autonomous energy-aware mode selection}, leveraging the efficiency advantage of wheeled locomotion on flat surfaces while preserving legged capability for complex terrain, directly addressing the operational limitation documented in deployed inspection systems \citep{Hutter2017ANYmal,Tranzatto2022CERBERUS}.

The remainder of this article is organized as follows: Section 2 details the methodology, including the physical robot design, the multi-body simulation setup, and the formal mathematical framework of the BG-inspired neural controller. Section 3 presents the experimental results, separating virtual multi-body simulation tests from physical prototype validations. Finally, Section 4 outlines the conclusions and directions for future work.

